\section*{摘要}\addcontentsline{toc}{section}{摘要}
本文比较HAR、GARCH(1,1)、XGB-Market、仅含截距及滞后残差的XGB-AR，以及加入MIDAS-EPU的动态误差修正模型DEC，预测下一交易日平方收益率的对数代理。公开修订版以解析梯度、多起点剖面最小二乘替代旧版单起点数值差分求解，并在特征构造前剔除上证50已确认的休市日记录。模型规格仍按验证集MSE选择，不以测试损失调整求解规则。

沪深300中，DEC的MSE为5.112644，高于XGB-Market的5.099716；QLIKE为3.536964，低于其4.951636。相对XGB-AR，三个市场的DEC均降低QLIKE、提高MSE；中国两市场的MAE下降，纳斯达克MAE上升。因此加入EPU提供的是本次设定下依赖损失函数的增量预测表现，不是全面优势或政策因果效应。纳斯达克的验证集最优规格由旧版12个月Beta改为12个月指数Almon；相对旧DEC，测试MSE下降而QLIKE上升。

重新计算五模型的30组滚动残差区间设置。共同1859个日期上，90\%区间、窗口250时，DEC得分10.1995，高于HAR的9.6973、XGB-Market的9.8534及XGB-AR的9.8793；DEC没有改善这些基准的区间质量。区间针对未来观测代理，不是潜在条件方差。

历史XGBoost的6888条预测已通过Windows列抽样兼容构建逐点重现；旧DEC仍不能逐点还原，历史环境不能唯一识别。新DEC在两套数值依赖下的最大预测差异小于$10^{-7}$，但多起点搜索不保证全局最优。由于这是已查看历史结果后的回顾性重分析，EPU发布时点、历史修订、完整交易日历及独立留出样本仍待核验；HAC及多重比较调整仅为探索性诊断。

\noindent\textbf{关键词：}波动率代理；XGBoost；MIDAS；经济政策不确定性；数值稳定性；消融比较
\clearpage
\section*{Abstract}\addcontentsline{toc}{section}{Abstract}
This revision compares HAR, GARCH(1,1), XGB-Market, a residual-only XGB-AR correction, and MIDAS-EPU dynamic error correction (DEC). The target is log next-day squared returns plus an offset. Analytic gradients and deterministic bounded multistart replace the legacy single-start numerical optimizer. A confirmed SSE50 holiday row is removed before constructing features and targets. Specification selection still uses validation MSE, without tuning the solver to test losses.

For CSI300, DEC has MSE 5.112644 and QLIKE 3.536964, compared with 5.099716 and 4.951636 for XGB-Market. Relative to XGB-AR, DEC lowers QLIKE but raises MSE in all three markets; MAE falls in the two Chinese samples and rises in Nasdaq. The EPU-augmented model therefore shows loss-dependent incremental performance under this design, not universal superiority or a causal policy effect. Nasdaq's selected specification changes from 12-month Beta to 12-month exponential Almon. Compared with legacy DEC, test MSE improves while QLIKE worsens.

All 30 rolling-residual interval settings are recomputed for five models. On 1,859 common dates, with 90\% intervals and a calibration window of 250, DEC scores 10.1995 versus 9.6973 for HAR, 9.8534 for XGB-Market and 9.8793 for XGB-AR. These intervals concern an observed proxy, not latent variance.

A Windows-compatible column sampler reproduces all 6,888 archived XGBoost forecasts exactly; the original DEC environment remains unidentified. Revised DEC forecasts differ by less than $10^{-7}$ across two numerical stacks. Multistart does not certify a global optimum. This retrospective reanalysis does not supply a fresh holdout or verified real-time EPU vintages. HAC and multiplicity-adjusted results remain exploratory diagnostics.
\clearpage
